\section{Appendix}
\subsection{Typumwandlung: \texttt{dict()} und \texttt{list()}}
\begin{itemize}
    \item \textbf{Tuple zu Dict}: Die Funktion \texttt{dict()} kann eine Liste von 2-Element-Tupeln direkt in Schlüssel-Wert-Paare umwandeln.
    \item \textbf{Duplikatentfernung}: Da Wörterbuchschlüssel eindeutig sein müssen, entfernt das Umwandeln einer Tupelliste in ein \texttt{dict} automatisch doppelte Einträge (wobei der zuletzt ausgewertete Wert für einen gegebenen Schlüssel beibehalten wird).
    \item \textbf{Dict zu List}: Wandelt das deduplizierte Wörterbuch zurück in eine Liste von Tupeln, indem \texttt{list(dictionary.items())} verwendet wird.
\end{itemize}

\lstinputlisting[language=Python]{code/type_coversion.py}

\subsection{Decorators}
Decorators sind Funktionen, die eine andere Funktion als Argument annehmen, um ihr Verhalten zu erweitern, ohne den ursprünglichen Quellcode zu verändern. Sie werden durch das Symbol \texttt{@} gekennzeichnet.

\begin{lstlisting}[language=Python]
def require_login(func):
    def wrapper(user):
        if user["logged_in"]:
            return func(user)
        print("Access denied.")
    return wrapper

@require_login
def view_dashboard(user):
    print(f"Welcome {user['name']}")
\end{lstlisting}


\subsection{Script Execution}
Um zu unterscheiden, ob ein Script direkt ausgeführt wird oder als Modul importiert wird, verwendet Python die Variable \texttt{\_\_name\_\_}.

\begin{itemize}
    \item Wenn direkt ausgeführt: \texttt{\_\_name\_\_ == "\_\_main\_\_"}.
    \item Wenn importiert: \texttt{\_\_name\_\_} entspricht dem Dateinamen des Moduls.
\end{itemize}

\begin{lstlisting}[language=Python]
def main():
    # Test code or entry point
    pass

if __name__ == "__main__":
    main()
\end{lstlisting}

\subsection{Distribuzioni Casuali in NumPy}

\begin{itemize}
    \item \textbf{np.random.rand}: Erzeugt Stichproben aus einer \textbf{Gleichverteilung} im Intervall $[0, 1)$. Jeder Wert hat die gleiche Wahrscheinlichkeit, gezogen zu werden.
    \item \textbf{np.random.randn}: Erzeugt Stichproben aus der \textbf{Standardnormalverteilung} (Gaußverteilung) mit dem Mittelwert $\mu = 0$ und der Varianz $\sigma^2 = 1$.
\end{itemize}

\subsubsection{Transformation der Normalverteilung}
Um die Standardnormalverteilung $\mathcal{N}(0, 1)$ zu einer benutzerdefinierten Verteilung $\mathcal{N}(\mu, \sigma^2)$ zu skalieren und zu verschieben, wird die folgende lineare Formel angewendet:

\[ X = \sigma \cdot \text{np.random.randn}(\dots) + \mu \]

Wobei:
\begin{itemize}
    \item $\sigma$ (Sigma): Standardabweichung (prüft die "Breite" der Glocke).
    \item $\mu$ (Mu): Erwartungswert (verschiebt das Zentrum der Verteilung).
\end{itemize}

\subsection{Indexing vs Slicing: Regola Rapida}

Der Unterschied hängt vom Typ des Selektors ab, der für die Achse verwendet wird:

\begin{itemize}
    \item \textbf{Ganzzahl (Indexing)}: Wenn du eine einzelne Zahl verwendest (z.B. \texttt{arr[2]}), wird die Dimension \textbf{entfernt} (rank-1).
    \item \textbf{Bereich (Slicing)}: Wenn du die beiden Doppelpunkte verwendest (z.B. \texttt{arr[2:3]}), wird die Dimension \textbf{beibehalten} (rank-2).
\end{itemize}

\subsubsection{Esempio Pratico}
Gegeben sei ein 2D-Array mit Shape \texttt{(3, 3)}:
\begin{itemize}
    \item \texttt{arr[2, :]} $\rightarrow$ Ergebnis 1D, Shape \texttt{(3,)}.
    \item \texttt{arr[2:3, :]} $\rightarrow$ Ergebnis 2D, Shape \texttt{(1, 3)}.
\end{itemize}